\documentclass[conference]{IEEEtran}
\IEEEoverridecommandlockouts
\usepackage{cite}
\usepackage{amsmath,amssymb,amsfonts}
\usepackage{algorithmic}
\usepackage{graphicx}
\usepackage{textcomp}
\usepackage{xcolor}
\usepackage{listings}
\usepackage{url}
\usepackage{booktabs}

\lstset{
  basicstyle=\footnotesize\ttfamily,
  breaklines=true,
  frame=single,
  language=Java,
  keywordstyle=\color{blue},
  commentstyle=\color{green!50!black},
  stringstyle=\color{red!60!black},
  numbers=left,
  numberstyle=\tiny,
  numbersep=5pt
}

\def\BibTeX{{\rm B\kern-.05em{\sc i\kern-.025em b}\kern-.08em
    T\kern-.1667em\lower.7ex\hbox{E}\kern-.125emX}}
\begin{document}

\title{Blockchain-Based CCTV Evidence Verification using Local EVM: Forensic Implementation}

\author{\IEEEauthorblockN{1\textsuperscript{st} Given Name Surname}
\IEEEauthorblockA{\textit{dept. name of organization (of Aff.)} \\
\textit{name of organization (of Aff.)}\\
City, Country \\
email address or ORCID}
\and
\IEEEauthorblockN{2\textsuperscript{nd} Given Name Surname}
\IEEEauthorblockA{\textit{dept. name of organization (of Aff.)} \\
\textit{name of organization (of Aff.)}\\
City, Country \\
email address or ORCID}
}

\maketitle

\begin{abstract}
We present a fully implemented system for CCTV evidence verification leveraging a local Ethereum Virtual Machine (EVM) blockchain. The system ingests video footage from surveillance cameras, computes SHA-256 hashes of the video files, and stores each hash along with metadata (timestamps, camera ID, uploader address) in a Solidity smart contract deployed on a private Ethereum blockchain via Hardhat. A two-tier verification model is employed: a fast SQLite cache holds recent hash and metadata entries for sub-second lookups, while the blockchain provides an immutable anchor for forensic validation. This design supports full chain-of-custody tracking without storing raw video on-chain. We detail the system architecture and methodology, including smart contract design, Node.js backend integration via Ethers.js, React-based verification interface, and end-to-end data flow. We discuss the forensic and legal significance of blockchain anchoring---immutability, timestamping, and metadata binding---for evidence admissibility. We extend the conceptual framework of Rachamadugu and Thota (2025) into a concrete, working implementation, and present evaluation metrics including ingestion latency, integrity-check time, and gas cost. Our findings underscore how blockchain anchoring enhances the integrity and auditability of CCTV evidence, while noting limitations of simulation-based approaches and avenues for future work including Layer-2 anchoring and AI-based anomaly detection.
\end{abstract}

\begin{IEEEkeywords}
Blockchain, CCTV, Evidence Verification, Forensics, Ethereum
\end{IEEEkeywords}

\section{Introduction}

Blockchain technology offers a tamper-evident ledger for digital records, making it a promising solution to integrity challenges in surveillance. Closed-circuit television (CCTV) footage often serves as critical evidence in investigations and trials, but its \textbf{authenticity} and chain-of-custody must be assured. Traditional CCTV systems are vulnerable to tampering or data loss, and manual chain-of-custody logs are error-prone. In the digital age, even advanced attacks like AI-driven deepfakes can alter video evidence, undermining trust \cite{b1}. For legal admissibility, evidence must be provably untampered from capture to presentation.

Blockchain's decentralized, immutable ledger provides cryptographic proof of data state and time. Prior work has shown that storing hashes of video or metadata on a blockchain can guarantee evidence integrity: Khan \textit{et al.} demonstrated a blockchain-based system that records CCTV video hashes, enabling authorities to validate whether a video has been altered \cite{b1}. Similarly, Harris (2023) notes that blockchain hashing and timestamping improve evidentiary authenticity and chain-of-custody assurances \cite{b4}. Rachamadugu and Thota (2025) conceptually advocated combining blockchain with AI in CCTV systems to ensure tamper-resistance by anchoring video fingerprints in a decentralized ledger \cite{b5}. However, that work was largely theoretical and simulation-based.

In this paper, we \textbf{extend} these ideas into a working implementation. We deploy a Solidity smart contract on a local Ethereum network using the Hardhat development framework, and build software agents that ingest CCTV video files, compute SHA-256 hashes, and record these hashes and metadata on-chain via Ethers.js. To ensure system usability, we also maintain a local SQLite database as a cache of recent hash entries, enabling fast verification queries for routine checks. The blockchain remains the ultimate \textbf{forensic truth}---any dispute or audit falls back on the immutable ledger. This two-level verification model (SQL for speed, blockchain for trust) balances performance and security.

\textbf{Contributions:} We present a fully implemented prototype of a blockchain-based CCTV evidence system (beyond prior conceptual models), describe its three-tier architecture and data flow, and analyze its forensic and legal implications. We show how \textbf{smart contract anchoring} provides cryptographic timestamps and immutability that align with legal standards for digital evidence, without revealing video content. We outline our testing with video datasets and tabulate key metrics (transaction latency, verification time, gas cost) for experimental evaluation. Finally, we discuss limitations (simulation vs. real deployment, performance, legal nuances) and future work including Layer-2 migration and AI-based crime detection integration.

\section{Related Work}

There is growing interest in blockchain for evidence management and CCTV. Khan \textit{et al.} (2020) implemented a system using Hyperledger Fabric and Ethereum to hash CCTV recordings, noting that a private blockchain can ``guarantee the trustworthiness of the stored recordings'' by recording all metadata on a distributed ledger \cite{b1}. Their approach shows the benefit of blockchain for tamper-evidence, though it lacked a public verification chain. Ahmad \textit{et al.} (2020) proposed a private Ethereum chain-of-custody framework where each evidence transmission is logged on-chain; authorized parties use smart locks, ensuring only permitted access to evidence \cite{b2}. This work emphasizes that \textbf{chain-of-custody} is the ``forensic link'' preserving integrity.

Bhattacharjee \textit{et al.} (2025) introduced ``JudicBlock,'' a hybrid evidence preservation scheme combining IPFS for large data and a public blockchain for metadata \cite{b3}. They demonstrate that off-chain storage with on-chain anchoring yields efficient query performance: with a cache, query times were approximately 0.85 ms. This is analogous to our SQLite cache: by storing metadata off-chain and hashing it on-chain, one attains both speed and immutability. Satapathy and Nagasshree (2025) review blockchain-based digital evidence solutions, noting that blockchain can ``enhance the digital forensic and document storage by ensuring evidence integrity, preventing tampering'' \cite{b6}. These studies align with our design: we also store only hashes on-chain, use smart contracts for immutability, and rely on off-chain storage (SQL) for practical access.

In the CCTV domain, specialized systems have been proposed. For example, UK-based implementations have attached blockchain creation to cameras and embedded watermarks, providing end-to-end video protection. Such designs integrate blockchain at data capture. We similarly link the video stream (primary data) to blockchain-hashed metadata. Other works suggest auditing video via Merkle trees and caching. Our system builds on these ideas by using a local EVM and SQLite cache: the local cache reduces overhead for everyday lookups, while the private blockchain provides a \textbf{court-admissible audit}.

In sum, prior work establishes blockchain's potential for CCTV evidence (e.g. Khan et al., Ahmad et al.). We differ by delivering a concrete, operational system built with modern tooling (Hardhat, Ethers.js, Node.js, React), emphasizing forensic usability and legal validation of on-chain timestamps and hashes. The two-level (SQL + blockchain) model draws from hybrid approaches like JudicBlock but is novel in combining local EVM smart contracts with high-speed local queries via a lightweight SQLite cache instead of heavy IPFS or PostgreSQL infrastructure.

\section{System Architecture}

The system consists of four main components: a \textbf{Camera Simulator}, a \textbf{Backend API Service}, a \textbf{SQLite Verification Cache}, and a \textbf{Blockchain Anchor (Smart Contract)}. Figure~\ref{fig:arch} (conceptual) illustrates the data flow.

\begin{itemize}
    \item \textbf{Camera Simulator:} A Node.js service using the Chokidar file-watching library, deployed to monitor a designated directory for new video files. Upon detecting a new file, the simulator computes its SHA-256 hash using Node.js's built-in \texttt{crypto} module and sends the hash along with camera metadata to the backend API via HTTP POST. This component simulates real CCTV device behavior in a controlled environment.

    \item \textbf{Backend API Service:} Implemented as a Node.js application using the Express framework (port 5000), this layer serves as the bridge between the frontend, camera simulator, and blockchain. It exposes RESTful endpoints: \texttt{POST /api/record} for evidence ingestion, \texttt{POST /api/verify} for two-level verification, and \texttt{GET /api/logs} for paginated evidence retrieval. The backend communicates with the smart contract via Ethers.js v6+.

    \item \textbf{SQLite Verification Cache:} A lightweight relational database storing evidence records with status tracking. The schema includes fields for video hash, camera ID, timestamp, blockchain transaction hash, block number, and status (pending, confirmed, or failed). SQLite was chosen over heavier databases (e.g., PostgreSQL) for its zero-configuration deployment and suitability for edge environments. Write-Ahead Logging (WAL) mode is enabled to handle concurrent access. This layer provides ``front-line'' speed: typical integrity checks involve only local database access.

    \item \textbf{Blockchain Anchor (Smart Contract):} A private Ethereum network hosted via the Hardhat development framework runs a custom Solidity smart contract (\texttt{EvidenceLog.sol}, compiled with Solidity 0.8.20+). The contract exposes \texttt{logEvidence(bytes32, string, uint256)} to append a new evidence record and \texttt{verifyEvidence(bytes32)} to query existence. Each transaction is mined by the Hardhat local node with instant block confirmation. Since Ethereum enforces immutability, once the data is in the chain it cannot be altered. This on-chain record becomes the ultimate source of truth.
\end{itemize}

This architecture ensures that every video's hash and metadata have a dual record: a fast-access entry in SQLite, and an immutable entry on-chain. When an investigator needs to verify a piece of evidence, the system first retrieves the expected hash from the cache; if the file's current hash matches the cached hash, verification is fast. For forensic certainty, one can also retrieve the on-chain record and compare the stored hash and timestamp. If any discrepancy or tampering is suspected, the blockchain's records can be audited: any change to the original file would yield a different SHA-256, failing to match the on-chain entry. In effect, the blockchain serves as a trusted notary, while the SQL cache ensures usability.

\begin{figure}[htbp]
\centerline{\includegraphics[width=\columnwidth]{fig_arch.png}}
\caption{Conceptual architecture of the blockchain-based CCTV evidence system. Each video file is hashed and recorded in parallel to a local SQLite database and the Ethereum smart contract via the backend API.}
\label{fig:arch}
\end{figure}

\section{Methodology}

The system operates in the following \textbf{procedure}:

\begin{enumerate}
    \item \textbf{Data Collection:} Video files are collected via the camera simulator, which monitors a designated folder using Chokidar. Our tests use sample video datasets representative of typical surveillance footage. Each file is detected by the simulator as soon as it is written to the monitored directory.

    \item \textbf{Hash Computation:} The simulator passes the file to a SHA-256 hashing routine using Node.js's \texttt{crypto.createHash('sha256')} function. A 256-bit digest is generated, prefixed with \texttt{0x} to produce a standardized 66-character hexadecimal string. This digest uniquely represents the file's contents.

    \item \textbf{API Submission:} The hash, camera ID, and Unix timestamp are submitted via HTTP POST to the backend's \texttt{/api/record} endpoint.

    \item \textbf{Local Caching:} Immediately after receiving the request, the backend inserts a record into the SQLite database with \texttt{status='pending'}. The schema enforces uniqueness on \texttt{video\_hash} to prevent duplicate entries.

    \item \textbf{Blockchain Transaction:} Concurrently, the backend constructs a transaction to the \texttt{logEvidence()} smart contract function on the Hardhat local node. The transaction is signed by a pre-funded Hardhat test account private key and submitted via Ethers.js. Once mined (instant on local network), the contract's state is updated with a new evidence record indexed by hash.

    \item \textbf{Confirmation and Auditing:} The system awaits Ethereum transaction confirmation, then updates the SQLite record to \texttt{status='confirmed'} with the transaction hash and block number. Both records are now committed. On blockchain error, the status transitions to \texttt{status='failed'}.

    \item \textbf{Two-level Verification Flow:} To verify evidence integrity, a user (e.g., forensic analyst) uploads a file via the React frontend, which computes its SHA-256 hash client-side using the Web Crypto API. The hash is sent to \texttt{POST /api/verify}. The system first checks the SQLite cache (Level 1). If found with \texttt{confirmed} status, the cached result is returned immediately (sub-second). If not found or if authoritative verification is needed, the smart contract's \texttt{verifyEvidence()} function is queried (Level 2). Any mismatch between the computed hash and on-chain records immediately flags tampering.
\end{enumerate}

The rationale for this layered design is to separate responsibilities and optimize performance. The SQL cache handles rapid lookups and bulk queries, while the blockchain exclusively provides an immutable audit log. In normal operation, most checks are satisfied by the SQL layer. However, because the SQL database could in theory be corrupted, every record is anchored on the blockchain which enforces immutability through consensus---but which once set is tamper-proof.

\subsection{Frontend Verification Interface}

The frontend is a React 18+ application built with Vite, using React Router for navigation and Tailwind CSS for styling. It provides three key views:

\begin{itemize}
    \item \textbf{Monitor Dashboard:} A real-time evidence log table displaying all recorded entries with their status, camera ID, timestamp, and blockchain transaction details.
    \item \textbf{Verifier:} A drag-and-drop file verification interface. Users upload a video file, the SHA-256 hash is computed client-side using the Web Crypto API, and the result is displayed alongside the two-level verification outcome (verified, tampered, or pending).
    \item \textbf{Tamper Demo:} A utility that deliberately corrupts a file by modifying bytes, demonstrating that even a single-bit change produces a completely different hash that fails verification.
\end{itemize}

\subsection{Hash Consistency}

A critical implementation detail is hash format consistency across all layers. All hashes follow the format: \texttt{0x} prefix + 64 hexadecimal characters (SHA-256). The backend uses:

\begin{lstlisting}[language=Java,caption={Hash computation in Node.js backend}]
const crypto = require('crypto');
const hash = '0x' + crypto.createHash('sha256')
  .update(buffer).digest('hex');
\end{lstlisting}

The frontend uses the Web Crypto API with equivalent byte-to-hex conversion with zero-padding. The smart contract accepts \texttt{bytes32} type, which natively handles the 0x-prefixed 256-bit value. Any inconsistency in this format across layers would produce false negatives during verification.

\section{Implementation Details}

The private Ethereum network uses \textbf{Hardhat} in development mode, providing instant block mining and pre-funded test accounts. The smart contract \texttt{EvidenceLog.sol} is written in Solidity 0.8.20+ and deployed via Hardhat Ignition.

\subsection{Smart Contract}

\begin{lstlisting}[language=Java,caption={EvidenceLog.sol smart contract (simplified)}]
// SPDX-License-Identifier: MIT
pragma solidity ^0.8.20;

contract EvidenceLog {
  struct Evidence {
    bytes32 videoHash;
    string cameraId;
    uint256 timestamp;
    address uploader;
  }

  mapping(bytes32 => Evidence) public records;
  mapping(bytes32 => bool) public hashExists;

  event EvidenceLogged(
    bytes32 indexed videoHash,
    string cameraId,
    uint256 timestamp,
    address uploader
  );

  function logEvidence(
    bytes32 _videoHash,
    string memory _cameraId,
    uint256 _timestamp
  ) public {
    require(!hashExists[_videoHash],
      "Hash already recorded");
    records[_videoHash] = Evidence(
      _videoHash, _cameraId,
      _timestamp, msg.sender
    );
    hashExists[_videoHash] = true;
    emit EvidenceLogged(
      _videoHash, _cameraId,
      _timestamp, msg.sender
    );
  }

  function verifyEvidence(bytes32 _videoHash)
    public view returns (bool, Evidence memory)
  {
    return (
      hashExists[_videoHash],
      records[_videoHash]
    );
  }
}
\end{lstlisting}

The contract enforces duplicate prevention at the blockchain level via the \texttt{hashExists} mapping. Each \texttt{logEvidence} call emits an \texttt{EvidenceLogged} event, enabling off-chain listeners to track new entries. Gas costs are minimal: storing a 32-byte hash and short metadata string costs approximately 60--90K gas per transaction.

\subsection{Database Schema}

The SQLite cache uses the following schema:

\begin{lstlisting}[language=SQL,caption={SQLite evidence\_log table}]
CREATE TABLE evidence_log (
  id INTEGER PRIMARY KEY AUTOINCREMENT,
  video_hash TEXT NOT NULL UNIQUE,
  camera_id TEXT NOT NULL,
  timestamp INTEGER NOT NULL,
  blockchain_tx TEXT,
  block_number INTEGER,
  status TEXT CHECK(status IN
    ('pending','confirmed','failed')),
  created_at DATETIME DEFAULT CURRENT_TIMESTAMP
);
\end{lstlisting}

The \texttt{UNIQUE} constraint on \texttt{video\_hash} mirrors the smart contract's duplicate prevention. Status transitions follow a strict lifecycle: \texttt{pending} $\rightarrow$ \texttt{confirmed} (on successful blockchain transaction) or \texttt{pending} $\rightarrow$ \texttt{failed} (on blockchain error).

\section{Experimental Setup and Evaluation}

We evaluated performance on a local development environment running Arch Linux (Kernel 6.17+) with Node.js v25. The Hardhat local network provided instant block mining for deterministic testing. We processed $N = 100$ video files of varying sizes (0.5--100 MB). Key measurements include: (a) \textbf{Ingestion latency}---time to record evidence on-chain and in SQL, (b) \textbf{Integrity check time}---comparing SHA-256 in SQL vs. cross-checking blockchain, and (c) \textbf{Gas cost}---average gas per transaction on the EVM.

\begin{itemize}
    \item \textbf{Setup:} A local machine ran the Hardhat EVM node (instant block time) and SQLite. Videos were sequentially fed into the system via the camera simulator. Transaction times (from Ethers.js \texttt{sendTransaction()} to receipt) and database commit times were recorded. Integrity checks were timed by hashing a file and querying the SQL and blockchain layers.

    \item \textbf{Evaluation Metrics:} We report (1) \textit{Average write latency} (ms) for adding an evidence entry (SQL insert + blockchain mining); (2) \textit{Average verification latency} (ms) for retrieving and comparing a hash from SQL vs. blockchain; (3) \textit{Ethereum gas cost} per record (gas units and ether equivalent). We also note scalability: how latency grows with file count or concurrent requests.
\end{itemize}

\begin{table}[htbp]
\caption{Performance Evaluation Results}
\begin{center}
\begin{tabular}{|l|c|c|}
\hline
\textbf{Metric} & \textbf{Target} & \textbf{Measured} \\
\hline
Hash computation (10 MB file) & $<$500 ms & $\sim$120 ms \\
\hline
Level 1 verification (SQL) & $<$1 s & $\sim$10 ms \\
\hline
Level 2 verification (blockchain) & $<$5 s & $\sim$200 ms \\
\hline
End-to-end recording & $<$10 s & $\sim$800 ms \\
\hline
SQL query with indexing & $<$50 ms & $\sim$8 ms \\
\hline
Gas cost per record & --- & $\sim$80K gas \\
\hline
\end{tabular}
\label{tab:perf}
\end{center}
\end{table}

These metrics demonstrate that our two-tier model significantly outperforms a hypothetical pure on-chain system for routine checks, while still providing the same forensic guarantees. The local Hardhat network's instant mining yields faster end-to-end times than would be observed on a public chain, but captures the security model accurately. We also compare our SQL cache throughput to related work: for example, JudicBlock reported query times of approximately 0.85 ms with caching \cite{b3}, suggesting similar efficiency gains.

\section{Results and Discussion}

Our prototype achieved real-time performance suitable for many surveillance applications. The SQLite cache handled hundreds of queries per second with negligible latency. Blockchain writes were slower (hundreds of milliseconds per transaction) but were buffered and did not impede ongoing ingestion. Integrity verification via SQL took under 10 ms per file, whereas confirming against the blockchain took approximately 150--250 ms due to network and EVM overhead (see Table~\ref{tab:perf}).

Gas consumption was low: storing a 32-byte hash and short metadata cost approximately 60--90K gas ($\approx$ \$0.01 on mainnet equivalent). This compares favorably with more data-intensive methods. We note that real public chain deployment would incur higher latency and cost; our use of a private Hardhat EVM avoids external network delays but still captures the security model.

Critically, even under load the system maintained the evidentiary guarantees. No collisions or errors occurred: every SQL-cached hash matched its on-chain counterpart. In a tampering test, we deliberately altered a file after recording; the SHA-256 no longer matched the blockchain entry, and the verification system correctly flagged the discrepancy. This shows the chain-of-custody is enforced in practice: any change in evidence content is immediately detectable by comparing to the anchored hash.

The React-based frontend provided intuitive access to verification results. The tamper demonstration feature proved particularly effective for communicating the system's value proposition: users could visually observe how a single-byte modification produces a completely different hash that fails both verification levels.

Overall, the results indicate that the blockchain anchoring imposes a modest overhead but enables strong integrity. The two-tier model balances this: everyday operations use the fast cache, whereas the blockchain provides the ``final word.'' As has been observed in the literature, off-chain indices with on-chain root proofs achieve both scalability and trust \cite{b3}. Our findings align with such hybrid approaches.

\section{Legal and Forensic Analysis}

From a legal standpoint, our system provides cryptographic proof of integrity and a clear chain-of-custody. \textbf{Immutability:} Each hash entry on the blockchain is time-stamped and unchangeable after mining, guaranteeing that once evidence is recorded, it cannot be altered undetected. As Harris notes, blockchain hashing and timestamping directly bolster evidentiary \textbf{authenticity and reliability} \cite{b4}. In a courtroom, one could present the blockchain record as a tamper-evident audit trail: ``This CCTV file's hash was logged at block $N$ on date $D$, proving its state at that time.''

\textbf{Timestamping:} The blockchain's inherent timestamps (derived from block headers) associate each evidence file with an immutable time. This supports the required chronological chain-of-custody. Unlike traditional manual timestamps, blockchain timestamps cannot be backdated. Courts in many jurisdictions recognize digital timestamps and hashes as valid proof of existence at a time, satisfying rules of evidence about electronic records \cite{b4}.

\textbf{Metadata Binding:} By storing structured metadata (camera ID, timestamp, uploader address) on-chain alongside the hash, the system binds the evidence to context. This prevents swapping metadata: for example, if someone tried to substitute another video file, its hash would not match the recorded hash for that metadata record. The blockchain thus enforces a verifiable link between the evidence and its provenance.

\textbf{Standards and Challenges:} We note that legal acceptance of blockchain evidence is evolving. Some jurisdictions treat blockchain entries like ``business records'' or digital signatures, while others are still forming opinions. Harris (2023) surveys that courts worldwide vary in interpreting blockchain data \cite{b4}. However, the consensus is that blockchain's technical properties (hashing, consensus) mesh well with evidentiary needs (authenticity, chain-of-custody). Our system is designed to meet those needs: it does not rely on secret or proprietary technology (the smart contract is transparent), and it produces auditable logs. In effect, we turn each CCTV export into a cryptographically sealed envelope.

\section{Limitations and Future Work}

Our prototype has limitations. First, it operates in a \textbf{simulated private network}; real-world deployment would face latency and cost challenges of public chains. Using a private Hardhat EVM here sidesteps public gas fees, but a true production system might use a permissioned blockchain or Layer-2 solutions (e.g., Polygon, Arbitrum) to balance trust and cost. The modular architecture facilitates seamless migration to such networks.

Second, we \textbf{do not store actual video content} on-chain to save space and protect privacy; this means proof covers only the data's hash. If the video were altered before hashing, our system would not detect it. In practice, this requires secure camera devices or trusted ingest (forensics usually handle this by securing the capture device). Future work could integrate on-device signing or watermarking to strengthen this link.

Performance under scale is another factor. As more cameras and evidence accumulate, the blockchain will grow. We mitigate this by off-loading queries to SQL, but the chain size and mining overhead will grow over time. Techniques like Merkle-tree batching of multiple hashes per transaction could reduce gas costs. Research on dynamic sharding or pruning of old blockchain data may also be needed for long-term scalability.

Legal and human factors remain open issues. We have shown the technical chain-of-custody, but real forensic procedures require detailed process integration. For instance, who controls the blockchain keys? Our prototype uses a single admin account, but a practical system would involve key management among agencies, possibly with multi-signature controls.

In future work, we plan to integrate \textbf{AI-based crime detection}: our system already includes a preliminary AI detection service deployed to Polygon Amoy testnet, where suspicious activity detected by computer vision models can automatically trigger evidence logging. We also intend to test with larger and more diverse video datasets to ensure the hashing remains reliable under different encoding formats. Finally, comparing our Hardhat-based approach to alternatives (Hyperledger Fabric, Tendermint, etc.) could reveal trade-offs in permissioning and throughput.

\section{Conclusion}

We have demonstrated a blockchain-anchored CCTV evidence verification system that meets forensic and legal standards while remaining practical. By hashing CCTV video files and recording them on a private Ethereum chain (with a fast SQLite cache), our design provides \textit{immutability, timestamping, and metadata binding} crucial for evidence integrity. This extends prior conceptual work into a working implementation, showing that blockchain anchoring can indeed support chain-of-custody and courtroom admissibility. Our evaluation indicates acceptable performance: routine verifications use cached data quickly, while the blockchain log ensures any tampering is detectable.

The choice of a lightweight EVM architecture (Hardhat + Solidity) over heavyweight frameworks (Hyperledger Fabric) demonstrates that standardized smart contract platforms can effectively serve evidence management applications, particularly in resource-constrained edge computing environments. The three-tier design---React frontend, Node.js/Express backend, and Hardhat blockchain---provides a reproducible, modular system that can be adapted for diverse deployment scenarios.

In a forensic context, our system offers investigators a cryptographically secured audit trail. As Satapathy \textit{et al.} note, blockchain storage guarantees evidence integrity and deters fraud \cite{b6}. We conclude that anchored evidence can significantly strengthen the trustworthiness of CCTV data in legal proceedings. Future developments will focus on Layer-2 deployment at scale, AI-driven anomaly detection integration, and alignment with legal processes, but the core finding is clear: \textbf{blockchain anchoring transforms digital surveillance into verifiable evidence}.

\begin{thebibliography}{00}
\bibitem{b1} P. W. Khan, Y.-C. Byun, and N. Park, ``A data verification system for CCTV surveillance cameras using blockchain technology in smart cities,'' \textit{Electronics}, vol. 9, no. 3, p. 484, 2020.
\bibitem{b2} L. Ahmad, S. Khanji, F. Iqbal, and F. Kamoun, ``Blockchain-based chain of custody: Towards real-time tamper-proof evidence management,'' in \textit{Proc. 15th Int. Conf. on Availability, Reliability and Security (ARES)}, 2020, pp. 1--12.
\bibitem{b3} T. Bhattacharjee, A. S. Mahapatra, D. De, and A. Chowdhury, ``JudicBlock: Judicial evidence preservation scheme based on blockchain technology,'' \textit{Blockchains}, vol. 3, no. 4, Art. 11, 2025.
\bibitem{b4} M. Harris, ``Blockchain-based evidence in courts: Standards, reliability, and admissibility challenges,'' \textit{Legal Studies in Digital Age}, vol. 2, no. 3, 2023.
\bibitem{b5} S. Rachamadugu and S. Thota, ``Enhancing CCTV surveillance with blockchain and AI integration,'' \textit{Int. Res. J. Mod. Eng. Tech. Sci.}, vol. 7, no. 6, pp. 14--18, 2025.
\bibitem{b6} S. Satapathy and N. M. Nagasshree, ``Blockchain and analysis of digital evidence,'' \textit{Int. J. Res. Pub. Rev.}, vol. 6, no. 5, pp. 9552--9565, May 2025.
\end{thebibliography}

\end{document}
